\section{Comparison clocks, scores and uncertainty}
\label{sup:protocol}\label{sup:protocol:availability}
\subsection{What is compared, and on which time origin}
The 2008 Fickian comparison starts at the measured critical-loading
crossing. The 2019 calculation retains the reported experimental clock,
with no time shift to reduce its falling-rate discrepancy; its initial
state and observable definitions are in Sec.~\ref{sup:fanerinputs}, and
its boundary calibration in Sec.~\ref{sup:fanermultiplier}. Crossing times
are interpolated linearly between successive states, and figure lines
interpolate the saved states. Main Table~\ref*{M-tab:journalmarch} gives
the measured and calculated crossings and falling interval. Numerical
controls and refinement evidence are in Secs.~\ref{sup:numericalcontrols}
and~\ref{sup:fanerrefinement}.

\subsection{Scores and uncertainty conventions}
Reading resolution, alternative source interpretations (mass basis, shape,
probe position and gas flow), and regression intervals are reported
separately; source alternatives are not statistical samples.
The band around a digitized point is the reading half-width of this digitization, not an experimental standard deviation; a point \emph{in band} lies within that resolution. The RMS drying score is the root mean square of model minus measured loading over the scored window, divided by the measured loading drop over that window. Changing a class distribution, coefficient or window defines a new calculation, not a new reading of the same score. The regression of converted uptake coefficients states 95-percent half-widths on its exponents and a residual standard deviation in log coefficient; neither covers the thesis samples' geometry or the \SI{41}{\kelvin} transfer. Faner's image-analysis counts give neither a mass-weighted histogram nor the thermocouple's local wetness, and the digitization half-widths represent neither uncertainty.

\section{Property domains and process boundaries}
\label{sup:propertyledger}
Constitutive domains and the positive-water continuation are specified in
main Sec.~\ref*{M-sec:thermo:properties}. Pure bulk hexane is evaluated
separately from the binary pore and surface state, with positive retained
water active after shell formation. The source's unreported water loading
and the adopted regime-A access assumption are defined in
Sec.~\ref{sup:fanerinputs}.
